\documentclass{article}%
\usepackage[T1]{fontenc}%
\usepackage[utf8]{inputenc}%
\usepackage{lmodern}%
\usepackage{textcomp}%
\usepackage{lastpage}%
\usepackage{geometry}%
\geometry{margin=2.2cm}%
%
\usepackage{amsmath}%
\usepackage{amssymb}%
\usepackage{booktabs}%
\usepackage{longtable}%
\usepackage{array}%
\usepackage{hyperref}%
\hypersetup{hidelinks}%
\title{IME 2027 {-} Química: comparação com o Poliedro\\[4pt]\large Comparação com o gabarito Poliedro\thanks{Gerado com gpt-resolve: \url{https://github.com/lgabs/gpt-resolve}}}%
\date{29/10/2026}%
%
\begin{document}%
\normalsize%
\maketitle%
\section*{Resumo}\textbf{Nota média: 10,0/10} (25 de 25 itens com a resolução correta; cada item vale o mesmo dentro da questão).\par\medskip%
\begin{longtable}{@{}c c p{0.78\textwidth}@{}}\toprule Questão & Nota & Resumo \\ \midrule
1 & 10,0 & A resolução e o gabarito Poliedro concordam na questão 1 e ambos estão certos. As duas reações coincidem: MnO2 + 4HCl $\to$ MnCl2 + Cl2 + 2H2O e Cl2 + 2NaBr $\to$ 2NaCl + Br2. A massa coincide: 0,870 g contra 0,87 g. Refiz a conta com as massas atômicas da folha de dados e obtive 0,8703 g. O enunciado transcrito confere com o original (pág. 2). A resolução está correta e completa. \\
2 & 10,0 & A resolução está correta nos dois itens e concorda com o gabarito Poliedro. No item a, o ciclo tem a mesma identificação das etapas: isotermas T1<T3 e duas adiabáticas, percorrido no sentido horário. As coordenadas do esboço foram conferidas numericamente e são consistentes. No item b, a demonstração W2+W4=0 usa o mesmo argumento: Q=0 nas adiabáticas, $\Delta$U depende só de T, e as duas etapas ligam T1 e T3 em sentidos opostos. A única diferença é a convenção de sinal do trabalho (pelo gás na resolução, sobre o gás no gabarito), e a própria resolução declara isso. Não há erro em nenhum dos dois. \\
3 & 10,0 & A resolução está correta nos quatro itens e bate com o gabarito Poliedro. As ordens são 2 em SO2 e 0 em O3, a lei é v = k{[}SO2{]}$^2$, k = 1,888 L$\cdot$mol$^{-}$$^{1}$$\cdot$s$^{-}$$^{1}$ e Ea = R$\cdot$ln4/(7/3000) = 594R $\approx$ 4,93 kJ/mol. Refiz as contas em python e todas conferem. Só há observações menores de redação (especulação sobre o mecanismo) e a hipótese implícita sobre a unidade de 7,56, nenhuma das quais afeta a nota. \\
4 & 10,0 & A resolução está correta e bate com o gabarito do Poliedro nos quatro itens. Em a e b, os dois dão as duas ligações de hidrogênio que formam um anel de 8 átomos. Em d, i = 0,58, b = 0,1156 mol/kg, $\Delta$T = 0,343 K e Tf = 5,11 $^\circ$C; refiz a conta em Python e os valores coincidem. A única diferença está no item c. Pela mesma tabela de equilíbrio (0,16 mol de M e 0,42 mol de D), os dois chegam a Kc = 16,4/C0 = 16,4$\cdot$V. A resolução escolhe C0 = 1 mol/L, declara a hipótese e informa 16,4 L/mol. O Poliedro deixa a resposta em função de V, porque falta esse dado no enunciado. As duas respostas são equivalentes. Fica em aberto se a banca aceita um valor numérico que depende de uma hipótese declarada. \\
5 & 10,0 & A resolução está certa e concorda com o gabarito Poliedro. Os dois chegam a $\Delta$n\_g = 450$\cdot$16,6/(8,3$\cdot$300) = +3 e concluem que a substância é o n{-}heptano (C7H16 + 11 O2 $\to$ 7 CO2 + 8 H2O). Recalculei com Python as contas e o $\Delta$n de cada candidato (-0,5, +1,5, +3, +0,5), e o sinal está bem justificado. A resolução é mais completa que o gabarito e não tem erro. Restam só observações menores sobre hipóteses idealizadas do enunciado. \\
6 & 10,0 & Resolução e gabarito Poliedro concordam nos três itens. (a) Dão a mesma equação de esterificação, balanceada 1:1:1:1 e com todas as espécies líquidas. (b) Dão $\alpha$ = 0,40 (40\%). (c) Dão Kc = 8/3, que a resolução ainda arredonda para $\approx$ 2,7. Refiz as contas: 4,0 mol de ácido e 2,0 mol de etanol, quadro de equilíbrio 2,4/0,4/1,6/1,6, K = 2,67, ou 2,69 sem arredondar. Estão corretas e não há erro em nenhum dos dois. Só ficaram observações menores sobre arredondamento e interpretação do enunciado. \\
7 & 10,0 & A resolução e o gabarito do Poliedro coincidem nos quatro itens: $\alpha$ = 1, $\beta$ = 2, $\gamma$ = -1; k = 1,25$\times$10$^{-}$$^2$ L$\cdot$mol$^{-}$$^{1}$$\cdot$s$^{-}$$^{1}$; v = k{[}A{]}{[}B{]}$^2$/{[}C{]}; ordem global 2. Os dois supõem o mesmo erro de digitação no v0 do exp. 4 (1,28$\times$10$^{-}$$^3$ no lugar do 1,28$\times$10$^{-}$$^{6}$ impresso). Refiz o sistema em Python. Com o valor corrigido, a solução é exatamente (1, 2, -1), e o k obtido reproduz os exps. 1--3. Com o valor impresso, os expoentes não fazem sentido químico ($\approx$ -1,49; 1,00; 0,49). A resolução está correta e bem justificada, inclusive pelo mecanismo (estado estacionário de X). O único ponto em aberto é se a banca vai anular a questão ou aceitar a correção do enunciado. \\
8 & 10,0 & A resolução está certa e coincide com o gabarito Poliedro nos dois itens. a) Y é o meso{-}3,4{-}dimetil{-}hexano. Conferi pelo produto misto, a partir das coordenadas do desenho, que o zigue{-}zague com um CH3 em cunha e o outro tracejado dá (3R,4S) e tem centro de inversão. A projeção de Fischer também é meso. b) X é o 2{-}clorobutano. A justificativa de menor número de carbonos está completa: todos os R de até 4 C foram testados e só o sec{-}butil gera carbonos quirais. A transcrição do enunciado confere com o PDF. Não encontrei erro de química nem de representação. \\
9 & 10,0 & A resolução concorda integralmente com o gabarito do Poliedro. A crioscopia dá M = 5,12$\cdot$53/(1,28$\cdot$2) = 106 g/mol, o que aponta para C8H10 ou C7H6O. O Tollens negativo exclui o benzaldeído, e o critério de produto único de halogenação seleciona o p{-}xileno (1,4{-}dimetilbenzeno). Conferi as contas e a análise de simetria, e a resolução e o gabarito estão corretos. \\
10 & 10,0 & A resolução está correta e bate com o gabarito Poliedro. Os dois usam o relógio Th--Pb, t = t$\frac12$$\cdot$ln(1+R)/ln2, com ln(1,024) = 7ln2 - 3ln5 = 0,024 pelos dados da prova. Isso dá t/t$\frac12$ = 0,03463 e t $\approx$ 4,88$\times$10\^{}8 anos. A resolução apresenta o resultado arredondado como 4,9$\times$10\^{}8 anos e comenta que, com ln exato, sairia $\approx$ 4,8$\times$10\^{}8 anos. A diferença para o gabarito é só de arredondamento, sem erro de conta nem de modelo. \\
\bottomrule\end{longtable}%
\section*{Principais pontos de dúvida}%
\begin{itemize}%
\item%
\textbf{Q4} \textit{(statement)} No item c o enunciado não dá volume nem concentração inicial, então Kc = 16,4/C0 (ou 16,4$\cdot$V) não tem valor único. A resolução adota C0 = 1 mol/L e chega a 16,4 L/mol. O Poliedro deixa a resposta em função de V. Outra leitura possível seria usar a concentração do item d (0,10 mol/L), que dá Kc $\approx$ 164 L/mol. Fica para o professor decidir se a banca aceita um valor numérico baseado numa hipótese declarada ou se exige a expressão literal. Isso pode mudar a nota do item c.%
\item%
\textbf{Q7} \textit{(statement)} O enunciado original (pág. 4 do PDF) traz v0 = 1,28$\times$10$^{-}$$^{6}$ no exp. 4. Com esse valor, o sistema dá $\alpha$ $\approx$ -1,49, $\beta$ $\approx$ 1,00, $\gamma$ $\approx$ 0,49, expoentes não inteiros e sem sentido químico. A mantissa 1,28 = 1,25$\times$1,024 coincide com 1,25$\times$2$^{1}$$^{0}$, o que indica que o expoente correto seria 10$^{-}$$^3$. O Poliedro sugere a anulação da questão. Vale levar ao professor se a banca vai anular a questão ou aceitar a correção. Isso pode mudar a pontuação de todos os itens.%
\end{itemize}%
\section*{Questão 1}%
\href{https://poliedroresolve.sistemapoliedro.com.br/vestibulares/ime/2027/ime-2-fase-dia-3-28-10-2026-2-fase-dia-3/dissertativa-1-quimica-2-fase-dia-3-ime-2027}{Resolução do gabarito}\par\smallskip%
\begin{longtable}{@{}p{0.06\textwidth} p{0.33\textwidth} p{0.33\textwidth} p{0.1\textwidth} p{0.08\textwidth}@{}}\toprule Item & Resolução & Gabarito & Veredito & Certo \\ \midrule
Reações envolvidas & MnO2 + 4 HCl $\to$ MnCl2 + Cl2 + 2 H2O; Cl2 + 2 NaBr $\to$ 2 NaCl + Br2 & MnO2 + 4HCl $\to$ MnCl2 + Cl2 + 2H2O; Cl2 + 2NaBr $\to$ 2NaCl + Br2 & igual & ambos \\
Massa de MnO2 & $\approx$ 0,870 g (0,0100 mol) & 0,87 g & igual & ambos \\
\bottomrule\end{longtable}%
\begin{itemize}%
\item%
\textbf{Reações envolvidas)} As equações são idênticas e estão balanceadas em massa e em carga. O Mn passa de +4 a +2 e ganha 2e$^{-}$; o Cl passa de -1 a 0, com 2e$^{-}$ por Cl2. O Cl2 desloca o Br$^{-}$ porque é oxidante mais forte.%
\item%
\textbf{Massa de MnO2)} Refiz a conta com python3. M(NaBr) = 102,894 g/mol, então n(NaBr) = 2,06/102,894 = 0,020021 mol. n(MnO2) = n(Cl2) = n(NaBr)/2 = 0,010010 mol. M(MnO2) = 86,936 g/mol, então m = 0,8703 g. Os valores 0,870 g e 0,87 g diferem só no número de algarismos significativos.%
\end{itemize}%
\paragraph{Pontos de dúvida}%
\begin{itemize}%
\item%
\textit{(solution)} A resolução supõe reagentes puros, rendimento de 100\% e uso de todo o Cl2 na oxidação do brometo, e declara essas hipóteses. O enunciado pede a massa de MnO2, não a de pirolusita, então a pureza do minério não interfere na resposta.%
\item%
\textit{(solution)} A resolução dá 0,870 g (3 algarismos significativos) e o gabarito dá 0,87 g. A diferença é só de arredondamento.%
\end{itemize}%
\section*{Questão 2}%
\href{https://poliedroresolve.sistemapoliedro.com.br/vestibulares/ime/2027/ime-2-fase-dia-3-28-10-2026-2-fase-dia-3/dissertativa-2-quimica-2-fase-dia-3-ime-2027}{Resolução do gabarito}\par\smallskip%
\begin{longtable}{@{}p{0.06\textwidth} p{0.33\textwidth} p{0.33\textwidth} p{0.1\textwidth} p{0.08\textwidth}@{}}\toprule Item & Resolução & Gabarito & Veredito & Certo \\ \midrule
a & Ciclo de Carnot percorrido no sentido horário no diagrama P$\times$V: I é a compressão isotérmica em T1 (1$\to$2); II é a compressão adiabática (2$\to$3, com T3>T1); III é a expansão isotérmica em T3 (3$\to$4); IV é a expansão adiabática de volta a 1 (4$\to$1). O esboço tem as isotermas T1<T3 e adiabáticas mais íngremes. & Gráfico P$\times$V do ciclo: 1$\to$2 (I, isoterma T1), 2$\to$3 (II, adiabática), 3$\to$4 (III, isoterma T3), 4$\to$1 (IV, adiabática). & igual & ambos \\
b & Convenção W = trabalho feito pelo gás. Com Q=0 nas adiabáticas, W=-$\Delta$U e U=U(T), então W2=-{[}U(T3)-U(T1){]}=-W4 e W2+W4=0. Logo Wtot=W1+W3=RT1 ln(V2/V1)+RT3 ln(V4/V3). Com adiabáticas reversíveis, Wtot=R(T3-T1) ln(V1/V2)>0. & W2=nCv(T3-T1) e W4=nCv(T1-T3)=-W2, logo Wtot=W1+W3, com W1=-nRT1 ln(V2/V1) e W3=-nRT3 ln(V4/V3). & equivalente & ambos \\
\bottomrule\end{longtable}%
\begin{itemize}%
\item%
\textbf{a)} A sequência de estados, o tipo de cada etapa e a posição relativa das isotermas são os mesmos do gabarito. T3>T1 está certo, porque a compressão adiabática aquece o gás. O sentido horário também está certo: a compressão ocorre na isoterma fria e a expansão na quente. Recalculei as coordenadas do TikZ com $\gamma$=5/3 e T3/T1=1,5: V3=1,0887, P3=1,3778, V4=2,7217 e P4=0,5511. Os pontos caem exatamente nas adiabáticas desenhadas, e V4/V3=V1/V2=2,5. O esboço é coerente.%
\item%
\textbf{b)} A diferença é só de convenção de sinal. O gabarito usa W = trabalho realizado sobre o gás ($\Delta$U=Q+W): W2=+Cv$\Delta$T>0 na compressão e W1=-RT1 ln(V2/V1)>0. A resolução usa W = trabalho realizado pelo gás ($\Delta$U=Q-W), e por isso todas as expressões aparecem com sinal trocado. A resolução declara isso explicitamente no fim. O ponto central coincide: Q=0 dá W=$\Delta$U (ou -$\Delta$U), e as duas adiabáticas ligam T1 e T3 em sentidos opostos, o que leva a W2+W4=0. A resolução é até mais geral, porque usa apenas U=U(T), sem exigir Cv constante, e vale também para adiabáticas irreversíveis. O resultado extra, R(T3-T1)ln(V1/V2), está corretamente condicionado a adiabáticas reversíveis (V4/V3=V1/V2 conferido numericamente).%
\end{itemize}%
\paragraph{Pontos de dúvida}%
\begin{itemize}%
\item%
\textit{(solution)} Convenção de sinal do trabalho: a resolução usa W realizado pelo gás (W1<0, W3>0), e o gabarito usa W realizado sobre o gás (W1=-RT1 ln(V2/V1)>0). A resolução avisa que na convenção oposta os sinais se invertem. Mesmo assim, um corretor que exija a forma do gabarito pode estranhar os sinais.%
\item%
\textit{(statement)} O enunciado chama as etapas II e IV de 'muito rápidas', o que sugere processos adiabáticos irreversíveis. Nesse caso não valem PV\^{}$\gamma$=cte nem V4/V3=V1/V2, e as curvas do gráfico seriam só indicativas. A resolução observa isso e mostra que o item b não depende da reversibilidade. O gabarito trata as etapas implicitamente como adiabáticas comuns.%
\item%
\textit{(key)} O gabarito escreve W2=nCv(T3-T1), o que pressupõe Cv constante. A demonstração da resolução usa apenas U=U(T) (lei de Joule), que é mais geral. As duas abordagens são aceitáveis para gás ideal.%
\end{itemize}%
\section*{Questão 3}%
\href{https://poliedroresolve.sistemapoliedro.com.br/vestibulares/ime/2027/ime-2-fase-dia-3-28-10-2026-2-fase-dia-3/dissertativa-3-quimica-2-fase-dia-3-ime-2027}{Resolução do gabarito}\par\smallskip%
\begin{longtable}{@{}p{0.06\textwidth} p{0.33\textwidth} p{0.33\textwidth} p{0.1\textwidth} p{0.08\textwidth}@{}}\toprule Item & Resolução & Gabarito & Veredito & Certo \\ \midrule
a & Ordem 2 em SO2 e ordem 0 em O3 (ordem global 2) & Ordem 2 em SO2, ordem 0 em O3 & igual & ambos \\
b & v = k{[}SO2{]}\^{}2 & v = k{[}SO2{]}\^{}2 & igual & ambos \\
c & k = 1,888 L$\cdot$mol$^{-}$$^{1}$$\cdot$s$^{-}$$^{1}$ (a 300 K) & k = 1,888 L$\cdot$mol$^{-}$$^{1}$$\cdot$s$^{-}$$^{1}$ & igual & ambos \\
d & Ea = R$\cdot$ln4/(1/300 - 1/1000) $\approx$ 4,93$\times$10$^3$ J/mol $\approx$ 4,9 kJ/mol ($\approx$ 1,18 kcal/mol) & Ea = 594R $\approx$ 4,93 kJ/mol & igual & ambos \\
\bottomrule\end{longtable}%
\begin{itemize}%
\item%
\textbf{a)} Nos experimentos 1 e 2, {[}O3{]} cai à metade e v não muda, então n = 0. Nos experimentos 2 e 3, {[}SO2{]} triplica e v fica 1,062/0,118 = 9 vezes maior, então m = 2. Conferi a tabela no PDF original: os dados batem com a transcrição.%
\item%
\textbf{b)} É a mesma lei de velocidade, consequência direta do item a.%
\item%
\textbf{c)} 0,118/0,25$^2$ = 1,888 e 1,062/0,75$^2$ = 1,888 (conferido em python). A unidade é a correta para ordem global 2.%
\item%
\textbf{d)} A razão é k1000/k300 = 7,56/1,888 = 4,004 $\approx$ 4, com ln 4 = 2$\cdot$ln 2 $\approx$ 1,386. Como 1/300 - 1/1000 = 7/3000 K$^{-}$$^{1}$, Ea/R = 1,386$\cdot$3000/7 $\approx$ 594 K. Isso dá Ea = 594$\cdot$8,3 $\approx$ 4931 J/mol $\approx$ 4,93 kJ/mol, o mesmo valor do gabarito (594R). Sem arredondar a razão para 4, o resultado seria 4935 J/mol, diferença desprezível. O valor alternativo de 1,18 kcal/mol (com R = 1,987) também confere.%
\end{itemize}%
\paragraph{Pontos de dúvida}%
\begin{itemize}%
\item%
\textit{(statement)} O enunciado dá a constante a 1000 K como '7,56 em módulo', sem unidade. A resolução supõe, sem prejuízo, que está na mesma unidade de k300 (L$\cdot$mol$^{-}$$^{1}$$\cdot$s$^{-}$$^{1}$). Como no cálculo só entra a razão k1000/k300, a hipótese é necessária, mas não afeta o resultado se as unidades forem as mesmas.%
\item%
\textit{(solution)} No item b, a afirmação de que 'o O3 só participa depois da etapa lenta, que é fotoquímica' e, na verificação, a leitura de que o Ea baixo vem de um processo ativado por UV são especulações de mecanismo que a questão não pede. Não estão erradas como hipótese, mas também não decorrem dos dados. O correto seria dizer apenas que a reação não é elementar.%
\end{itemize}%
\section*{Questão 4}%
\href{https://poliedroresolve.sistemapoliedro.com.br/vestibulares/ime/2027/ime-2-fase-dia-3-28-10-2026-2-fase-dia-3/dissertativa-4-quimica-2-fase-dia-3-ime-2027}{Resolução do gabarito}\par\smallskip%
\begin{longtable}{@{}p{0.06\textwidth} p{0.33\textwidth} p{0.33\textwidth} p{0.1\textwidth} p{0.08\textwidth}@{}}\toprule Item & Resolução & Gabarito & Veredito & Certo \\ \midrule
a & Ligações de hidrogênio O--H$\cdot$$\cdot$$\cdot$O=C: são duas e fecham um anel de 8 átomos. & Duas ligações de hidrogênio entre as carboxilas. & igual & ambos \\
b & Duas moléculas de CH3--C(=O)--O--H unidas em anel por duas ligações de hidrogênio O--H$\cdot$$\cdot$$\cdot$O=C, com 2 pares isolados em cada O e 48 elétrons de valência no total. & Dímero cíclico com duas ligações de H (estrutura em imagem). & igual & ambos \\
c & Kc = 0,42/(0,16)$^2$ $\approx$ 16,4 L/mol, tomando como base 1 mol/L. Na forma geral, Kc = 16,4/C0. & Kc = 0,42V/0,16$^2$ = 16,4$\cdot$V (V em L). Segundo o Poliedro, sem volume nem concentração não há valor numérico. & equivalente & ambos \\
d & i = 0,58; b $\approx$ 0,1156 mol/kg; $\Delta$Tf $\approx$ 0,343 K; Tf $\approx$ 5,11 $^\circ$C. & i = 0,58; W = 0,1156 mol/kg; $\Delta$Tc = 0,343 $^\circ$C; T = 5,11 $^\circ$C. & igual & ambos \\
\bottomrule\end{longtable}%
\begin{itemize}%
\item%
\textbf{a)} Mesma conclusão: no dímero cíclico há duas ligações de hidrogênio entre as carboxilas (H da hidroxila de uma molécula com o O carbonílico da outra). A resolução ainda explica por que o benzeno favorece o dímero, o que está correto.%
\item%
\textbf{b)} Conferi o desenho em TikZ. O anel é C1=Oa$\cdot$$\cdot$$\cdot$Hb--Ob--C2=Od$\cdot$$\cdot$$\cdot$Hc--Oc--C1, com 8 átomos. Cada C e cada O fecham o octeto e cada O tem 2 pares isolados. A conta de 24 e$^{-}$ por monômero (C: 8, H: 4, O: 12), 48 no dímero, está certa. É a mesma estrutura do gabarito.%
\item%
\textbf{c)} Os dois usam o mesmo quadro de equilíbrio: por mol de ácido, sobram 0,16 mol de M e formam{-}se 0,42 mol de D. Daí Kc = (0,42/V)/(0,16/V)$^2$ = 0,42V/0,0256 = 16,406$\cdot$V. A resolução escreve Kc = 16,4/C0. Com C0 = 1 mol/V, isso dá exatamente 16,4$\cdot$V, a expressão do gabarito. O valor 16,4 L/mol sai só da hipótese V = 1 L (C0 = 1 mol/L), que a resolução declara e justifica. O gabarito, por sua vez, prefere deixar Kc em função de V. Os dois reconhecem que o enunciado não traz esse dado. A diferença está só em como apresentam o resultado.%
\item%
\textbf{d)} Refiz a conta com M(CH3COOH) = 60,05 g/mol. Em 1 L de solução, 0,10 mol de ácido pesam 6,005 g e ocupam 5,72 mL. Sobram 994,28 mL de benzeno, ou 865,0 g. Então b = 0,10/0,8650 = 0,11560 mol/kg. Com i = 0,16 + 0,42 = 0,58, $\Delta$T = 5,12$\cdot$0,58$\cdot$0,1156 = 0,3433 K e Tf = 5,45 - 0,343 = 5,107 $\approx$ 5,11 $^\circ$C. Todos os valores intermediários e o final coincidem com o gabarito.%
\end{itemize}%
\paragraph{Pontos de dúvida}%
\begin{itemize}%
\item%
\textit{(statement)} No item c o enunciado não dá volume nem concentração inicial, então Kc = 16,4/C0 (ou 16,4$\cdot$V) não tem valor único. A resolução adota C0 = 1 mol/L e chega a 16,4 L/mol. O Poliedro deixa a resposta em função de V. Outra leitura possível seria usar a concentração do item d (0,10 mol/L), que dá Kc $\approx$ 164 L/mol. Fica para o professor decidir se a banca aceita um valor numérico baseado numa hipótese declarada ou se exige a expressão literal. Isso pode mudar a nota do item c.%
\item%
\textit{(statement)} Os itens c e d não são coerentes entre si. Com Kc constante, a razão dímero/monômero mudaria ao passar para 0,10 mol/L. Por isso o item d precisa impor a mesma razão (i = 0,58). Tanto a resolução quanto o Poliedro apontam isso, e o item d foi resolvido seguindo essa imposição.%
\item%
\textit{(statement)} A frase 'cada mol de monômero gera 0,42 mol de dímero' foi lida, pela resolução e pelo gabarito, como estequiometria a partir de 1 mol de ácido (0,84 mol consumidos, 0,16 mol de M e 0,42 mol de D no equilíbrio). Mas o item d fala em 'relação dímero/monômero', o que abre outra leitura: {[}D{]}/{[}M{]} = 0,42 no equilíbrio. Nessa leitura, i $\approx$ 0,772, Tf $\approx$ 4,99 $^\circ$C e Kc $\approx$ 0,773/C0. A leitura adotada é a mais natural para o verbo 'gerar' e é a mesma do gabarito.%
\item%
\textit{(solution)} Para passar de 0,10 mol/L a molalidade, o item d supõe volumes aditivos e usa a densidade do benzeno puro: 0,865 kg de benzeno por litro, b = 0,1156 mol/kg. Desprezando o volume do ácido, b = 0,1149 mol/kg e Tf $\approx$ 5,109 $^\circ$C. A resposta arredondada continua 5,11 $^\circ$C, a mesma do gabarito.%
\end{itemize}%
\section*{Questão 5}%
\href{https://poliedroresolve.sistemapoliedro.com.br/vestibulares/ime/2027/ime-2-fase-dia-3-28-10-2026-2-fase-dia-3/dissertativa-5-quimica-2-fase-dia-3-ime-2027}{Resolução do gabarito}\par\smallskip%
\begin{longtable}{@{}p{0.06\textwidth} p{0.33\textwidth} p{0.33\textwidth} p{0.1\textwidth} p{0.08\textwidth}@{}}\toprule Item & Resolução & Gabarito & Veredito & Certo \\ \midrule
Questão 5 (identificação da substância) & n{-}heptano (C7H16). $\Delta$H - $\Delta$U = 450$\cdot$16,6 = 7470 J = $\Delta$n$\cdot$RT $\Rightarrow$ $\Delta$n = 7470/(8,3$\cdot$300) = +3; C7H16 + 11 O2 $\to$ 7 CO2 + 8 H2O. Os demais candidatos dão -0,5 (acetileno), +1,5 (p{-}xileno) e +0,5 (propeno). & $\Delta$n\_gás = 450$\cdot$16,6/(8,3$\cdot$300) = 3 $\Rightarrow$ n{-}heptano (C7H16 + 11 O2 $\to$ 7 CO2 + 8 H2O) & igual & ambos \\
\bottomrule\end{longtable}%
\begin{itemize}%
\item%
\textbf{Questão 5 (identificação da substância))} A conclusão (n{-}heptano) e o raciocínio ($\Delta$n\_g = 3) são os mesmos nas duas. Refiz as contas: 450$\cdot$16,6 = 7470 J e 8,3$\cdot$300 = 2490 J/mol, então $\Delta$n = 3,0. Pela fórmula geral $\Delta$n = y/4 - 1 para CxHy com H2O gasosa, os valores são: acetileno -0,5, p{-}xileno +1,5, n{-}heptano +3,0 e propeno +0,5. Todas as equações balanceadas estão certas. O sinal também está certo: $\Delta$T\_V > $\Delta$T\_P significa |$\Delta$U| > |$\Delta$H|. Como $\Delta$H - $\Delta$U = $\Delta$nRT > 0, há expansão gasosa. A resolução justifica melhor que o gabarito: mostra os quatro candidatos e explica o sinal. Conferi o enunciado transcrito com o PDF original e ele está fiel.%
\end{itemize}%
\paragraph{Pontos de dúvida}%
\begin{itemize}%
\item%
\textit{(statement)} O enunciado considera n{-}heptano e p{-}xileno gasosos a 27 $^\circ$C (item ii), embora sejam líquidos nessas condições a 1 atm (os pontos de ebulição são cerca de 98 $^\circ$C e 138 $^\circ$C). Isso é uma hipótese do enunciado, que a resolução segue com razão, inclusive tratando a H2O como gasosa. Não altera a resposta.%
\item%
\textit{(solution)} Algumas idealizações ficaram implícitas: todo o calor da reação vai para o calorímetro, a capacidade térmica do conteúdo é desprezada, os gases são ideais e o termo $\Delta$nRT é calculado a T = 300 K, embora o calorímetro aqueça cerca de 10$^2$ K. São as hipóteses{-}padrão do problema e também estão implícitas no gabarito.%
\end{itemize}%
\section*{Questão 6}%
\href{https://poliedroresolve.sistemapoliedro.com.br/vestibulares/ime/2027/ime-2-fase-dia-3-28-10-2026-2-fase-dia-3/dissertativa-6-quimica-2-fase-dia-3-ime-2027}{Resolução do gabarito}\par\smallskip%
\begin{longtable}{@{}p{0.06\textwidth} p{0.33\textwidth} p{0.33\textwidth} p{0.1\textwidth} p{0.08\textwidth}@{}}\toprule Item & Resolução & Gabarito & Veredito & Certo \\ \midrule
a & CH3COOH(l) + C2H5OH(l) ? CH3COOC2H5(l) + H2O(l), coeficientes 1:1:1:1 & CH3COOH(l) + C2H5OH(l) ? CH3COOC2H5(l) + H2O(l) & igual & ambos \\
b & $\alpha$(ácido acético) = 1,6/4,0 = 0,40 (40\%) & $\alpha$(ácido acético) = 1,6/4 = 0,4 (40\%) & igual & ambos \\
c & Kc = 1,6$^2$/(2,4$\cdot$0,4) = 8/3 $\approx$ 2,7 (adimensional) & Kc = (1,6$\cdot$1,6)/(2,4$\cdot$0,4) = 8/3 & igual & ambos \\
\bottomrule\end{longtable}%
\begin{itemize}%
\item%
\textbf{a)} É a mesma equação de esterificação, com as mesmas fases (todas líquidas, como pede a hipótese i) e o mesmo balanceamento 1:1:1:1. Conferi os átomos: C4H10O3 de cada lado.%
\item%
\textbf{b)} Os números de mols batem: n(ácido) = 240/60,05 $\approx$ 3,997 $\approx$ 4,0 mol e n(etanol) = 92/46,07 $\approx$ 1,997 $\approx$ 2,0 mol. Logo $\alpha$ = 1,6/4,0 = 0,40. Sem arredondar, dá 0,4003, o que não muda nada.%
\item%
\textbf{c)} Os dois montam o mesmo quadro de equilíbrio: 2,4 / 0,4 / 1,6 / 1,6 mol. Os dois chegam a 8/3 $\approx$ 2,67, e o 2,7 da resolução é só arredondamento. Como $\Delta$n = 0, o volume se cancela. Com as quantidades sem arredondar, recalculei K $\approx$ 2,69. A água entra corretamente na expressão, porque é produto e não solvente em excesso.%
\end{itemize}%
\paragraph{Pontos de dúvida}%
\begin{itemize}%
\item%
\textit{(solution)} No item c, a resolução escreve '8/3 $\approx$ 2,7'. O valor exato é 2,67, e com as massas molares sem arredondar dá cerca de 2,69. Isso vem só de arredondamento, mas convém deixar 8/3 ($\approx$ 2,67) como resposta principal, igual ao gabarito.%
\item%
\textit{(statement)} O enunciado não define o volume nem a fase de cada espécie. A hipótese (i), 'somente equilíbrio líquido{-}líquido', leva a tratar as quatro espécies como líquidas e a pôr a água no Kc. Como $\Delta$n = 0, Kc não depende de V, e isso a resolução justifica de forma explícita.%
\item%
\textit{(statement)} O termo 'grau de equilíbrio do sistema reacional para o ácido acético' poderia ser lido em relação ao reagente limitante (o etanol, com 80\%). Mas o enunciado pede de forma explícita o valor para o ácido acético, que é 40\%. Resolução e gabarito concordam, e a resolução cita os 80\% do etanol só como comparação.%
\end{itemize}%
\section*{Questão 7}%
\href{https://poliedroresolve.sistemapoliedro.com.br/vestibulares/ime/2027/ime-2-fase-dia-3-28-10-2026-2-fase-dia-3/dissertativa-7-quimica-2-fase-dia-3-ime-2027}{Resolução do gabarito}\par\smallskip%
\begin{longtable}{@{}p{0.06\textwidth} p{0.33\textwidth} p{0.33\textwidth} p{0.1\textwidth} p{0.08\textwidth}@{}}\toprule Item & Resolução & Gabarito & Veredito & Certo \\ \midrule
a & $\alpha$ = 1, $\beta$ = 2, $\gamma$ = -1, supondo que o v0 do exp. 4 seja 1,28$\times$10$^{-}$$^3$ e não o 1,28$\times$10$^{-}$$^{6}$ impresso & $\alpha$ = 1, $\beta$ = 2, $\gamma$ = -1, supondo erro de digitação no v0 do exp. 4 (1,28$\cdot$10$^{-}$$^3$ em vez de 1,28$\cdot$10$^{-}$$^{6}$) & igual & ambos \\
b & k = 1,25$\times$10$^{-}$$^2$ L$\cdot$mol$^{-}$$^{1}$$\cdot$s$^{-}$$^{1}$ & k = 1,25$\cdot$10$^{-}$$^2$ L$\cdot$mol$^{-}$$^{1}$$\cdot$s$^{-}$$^{1}$ & igual & ambos \\
c & v\_P = 1,25$\times$10$^{-}$$^2$$\cdot${[}A{]}{[}B{]}$^2$/{[}C{]}, com k = k1k2/k3 & v = k{[}A{]}{[}B{]}$^2${[}C{]}$^{-}$$^{1}$ & igual & ambos \\
d & ordem global = 2 & ordem global = 2 & igual & ambos \\
\bottomrule\end{longtable}%
\begin{itemize}%
\item%
\textbf{a)} Refiz as contas. Dividindo pelo exp. 1, chega{-}se a: $\alpha$+2$\beta$+3$\gamma$ = 2; 2$\alpha$+$\beta$+4$\gamma$ = 0; 3$\alpha$+4$\beta$+$\gamma$ = log$_{2}$(v4/v1). Com v4 = 1,28$\times$10$^{-}$$^3$, a razão é 1024 = 2$^{1}$$^{0}$ e a solução exata é (1, 2, -1). Com o valor impresso, a razão é 1,024 e a solução vira $\alpha$ $\approx$ -1,49, $\beta$ $\approx$ 1,00, $\gamma$ $\approx$ 0,49, valores sem sentido químico. A resolução chega a esses mesmos números. Os dois adotam a mesma hipótese de erro de digitação, e a resolução ainda a justifica pelo mecanismo, com o estado estacionário de X no limite k3{[}C{]} ? k2{[}B{]}.%
\item%
\textbf{b)} Pelo exp. 1: k = 1,25$\times$10$^{-}$$^{6}$ / {[}(10$^{-}$$^2$)(10$^{-}$$^2$)$^2$/(10$^{-}$$^2$){]} = 1,25$\times$10$^{-}$$^2$. A unidade L$\cdot$mol$^{-}$$^{1}$$\cdot$s$^{-}$$^{1}$ é a de ordem global 2. Com esse k, a lei reproduz exatamente os exps. 2 e 3 (5,00$\times$10$^{-}$$^{6}$ e 1,25$\times$10$^{-}$$^{6}$) e prevê 1,28$\times$10$^{-}$$^3$ para o exp. 4.%
\item%
\textbf{c)} É a mesma lei. A resolução substitui o valor numérico de k e mostra de onde ele vem (k = k1k2/k3, pela aproximação do estado estacionário). O gabarito deixa k simbólico. A diferença é só de forma.%
\item%
\textbf{d)} 1 + 2 - 1 = 2. O resultado é coerente com a unidade de k.%
\end{itemize}%
\paragraph{Pontos de dúvida}%
\begin{itemize}%
\item%
\textit{(statement)} O enunciado original (pág. 4 do PDF) traz v0 = 1,28$\times$10$^{-}$$^{6}$ no exp. 4. Com esse valor, o sistema dá $\alpha$ $\approx$ -1,49, $\beta$ $\approx$ 1,00, $\gamma$ $\approx$ 0,49, expoentes não inteiros e sem sentido químico. A mantissa 1,28 = 1,25$\times$1,024 coincide com 1,25$\times$2$^{1}$$^{0}$, o que indica que o expoente correto seria 10$^{-}$$^3$. O Poliedro sugere a anulação da questão. Vale levar ao professor se a banca vai anular a questão ou aceitar a correção. Isso pode mudar a pontuação de todos os itens.%
\item%
\textit{(solution)} A resolução justifica os expoentes pelo limite k3{[}C{]} ? k2{[}B{]}. No exp. 4, porém, {[}B{]}/{[}C{]} = 8, então o limite exige k3 ? 8k2. É uma hipótese implícita, sem dado no enunciado. Ela não afeta o resultado porque os expoentes também saem das próprias razões experimentais (exps. 1--3, mais o exp. 4 corrigido).%
\item%
\textit{(key)} O texto de resolução do gabarito ('sugere anulação... sem solução viável sem calculadora') parece contradizer as respostas por item, que supõem o erro de digitação e chegam a (1, 2, -1). Provavelmente o Poliedro deu as respostas sob a hipótese do erro e, ao mesmo tempo, recomendou a anulação. Não há divergência de conteúdo com a resolução.%
\end{itemize}%
\section*{Questão 8}%
\href{https://poliedroresolve.sistemapoliedro.com.br/vestibulares/ime/2027/ime-2-fase-dia-3-28-10-2026-2-fase-dia-3/dissertativa-8-quimica-2-fase-dia-3-ime-2027}{Resolução do gabarito}\par\smallskip%
\begin{longtable}{@{}p{0.06\textwidth} p{0.33\textwidth} p{0.33\textwidth} p{0.1\textwidth} p{0.08\textwidth}@{}}\toprule Item & Resolução & Gabarito & Veredito & Certo \\ \midrule
a & Y = meso{-}3,4{-}dimetil{-}hexano, (3R,4S), C8H18. Desenhado em zigue{-}zague (CH3 de C3 em cunha, CH3 de C4 tracejado) e em projeção de Fischer (os dois CH3 do mesmo lado, com plano de simetria). & Y = forma meso do 3,4{-}dimetil{-}hexano (a estrutura 3D está em imagem) & igual & ambos \\
b & X = 2{-}clorobutano, CH3CHClCH2CH3 & X = 2{-}clorobutano & igual & ambos \\
\bottomrule\end{longtable}%
\begin{itemize}%
\item%
\textbf{a)} Pela síntese de Wurtz (2 R--Cl + 2 Na $\to$ R--R + 2 NaCl), Y é um alcano simétrico R--R. Para Y ter dois C quirais e não desviar a luz polarizada, precisa ser a forma meso. Com R de até 3 C (metil, etil, propil, isopropil), nenhum produto tem carbono quiral. Com R de 4 C, só o sec{-}butil dá centros quirais, gerando o 3,4{-}dimetil{-}hexano (C8). Então o menor Y é mesmo o meso{-}3,4{-}dimetil{-}hexano. Conferi o desenho da resolução calculando o sinal do produto misto a partir das coordenadas do TikZ. C3 dá o mesmo sinal da referência R e C4 dá o sinal oposto (S), ou seja, (3R,4S), que é meso. A inversão pelo ponto médio de C3--C4 leva o CH3 em cunha de C3 ao CH3 tracejado de C4 e C2 a C5, logo há centro de inversão. A projeção de Fischer, com os dois CH3 à esquerda e os dois H à direita, também dá (3R,4S) e tem plano de simetria horizontal. Coincide com o gabarito.%
\item%
\textbf{b)} Para que R--R seja o 3,4{-}dimetil{-}hexano, R tem de ser o sec{-}butil (butan{-}2{-}il). Logo X é o 2{-}clorobutano, como diz o gabarito.%
\end{itemize}%
\paragraph{Pontos de dúvida}%
\begin{itemize}%
\item%
\textit{(statement)} O enunciado pede um Y que 'possui dois carbonos quirais e não desvia a luz polarizada'. A leitura pretendida, seguida pela resolução e pelo gabarito, é a forma meso. Só que o Wurtz partindo de 2{-}clorobutano racêmico (e por via radicalar/organossódica) dá uma mistura de meso e ($\pm$){-}3,4{-}dimetil{-}hexano, e o racemato também é opticamente inativo. Isso não muda a resposta esperada. A resolução trata o assunto só de passagem, na observação sobre o acoplamento R + S.%
\item%
\textit{(key)} O texto do Poliedro disponível não traz a estrutura 3D nem o nome de Y (estão em imagem). Por isso não dá para conferir a representação espacial do gabarito, apenas que se trata da forma meso do 3,4{-}dimetil{-}hexano.%
\end{itemize}%
\section*{Questão 9}%
\href{https://poliedroresolve.sistemapoliedro.com.br/vestibulares/ime/2027/ime-2-fase-dia-3-28-10-2026-2-fase-dia-3/dissertativa-9-quimica-2-fase-dia-3-ime-2027}{Resolução do gabarito}\par\smallskip%
\begin{longtable}{@{}p{0.06\textwidth} p{0.33\textwidth} p{0.33\textwidth} p{0.1\textwidth} p{0.08\textwidth}@{}}\toprule Item & Resolução & Gabarito & Veredito & Certo \\ \midrule
Massa molar de A & M = 5,12$\cdot$53/(1,28$\cdot$2) = 106 g/mol & M = 106 g/mol & igual & ambos \\
Estrutura molecular de A & p{-}xileno (1,4{-}dimetilbenzeno), C8H10 (C6H4(CH3)2) & p{-}xileno (1,4{-}dimetilbenzeno) & igual & ambos \\
\bottomrule\end{longtable}%
\begin{itemize}%
\item%
\textbf{Massa molar de A)} Refiz a conta: M = K\_c$\cdot$m\_A/($\Delta$T$\cdot$m\_benzeno) = 5,12$\cdot$53/(1,28$\cdot$2) = 106,0 g/mol. Supondo i = 1 (soluto não eletrólito e sem associação em benzeno), a conta confere.%
\item%
\textbf{Estrutura molecular de A)} Com M = 106, as fórmulas candidatas são C8H10 (M = 106,17) e C7H6O (benzaldeído, M = 106,12). O Tollens negativo descarta o aldeído. Entre os isômeros de C8H10 (etilbenzeno e o{-}, m{-} e p{-}xileno), só o p{-}xileno tem os 4 H do anel equivalentes e os dois CH3 equivalentes, e por isso dá um único produto de monohalogenação, seja no anel (X2/FeX3) ou na cadeia lateral (X2/h$\nu$). O raciocínio está correto e completo.%
\end{itemize}%
\paragraph{Pontos de dúvida}%
\begin{itemize}%
\item%
\textit{(statement)} O enunciado diz apenas que há 'reação de halogenação gerando um único produto possível'. Não informa as condições (anel com catalisador ou cadeia lateral com luz) nem o grau de substituição. Entende{-}se que é monohalogenação. A resolução trata bem a ambiguidade: mostra que o p{-}xileno dá produto único nas duas condições.%
\item%
\textit{(solution)} A resolução considera só compostos benzênicos de C, H e O. Em tese, existem outras fórmulas aromáticas com M $\approx$ 106 contendo N (por exemplo, C6H6N2), que não foram discutidas. É a hipótese padrão desse tipo de questão e não afeta a resposta esperada.%
\item%
\textit{(solution)} A resolução supõe i = 1 (sem dissociação ou associação de A em benzeno) e usa as massas inteiras para chegar a 106. A massa exata do C8H10 pela tabela fornecida é 106,17 g/mol, compatível com o dado.%
\end{itemize}%
\section*{Questão 10}%
\href{https://poliedroresolve.sistemapoliedro.com.br/vestibulares/ime/2027/ime-2-fase-dia-3-28-10-2026-2-fase-dia-3/dissertativa-10-quimica-2-fase-dia-3-ime-2027}{Resolução do gabarito}\par\smallskip%
\begin{longtable}{@{}p{0.06\textwidth} p{0.33\textwidth} p{0.33\textwidth} p{0.1\textwidth} p{0.08\textwidth}@{}}\toprule Item & Resolução & Gabarito & Veredito & Certo \\ \midrule
Idade da amostra (relógio 208Pb/232Th) & t = (t$\frac12$/ln2)$\cdot$ln(1,024) = (1,41$\times$10\^{}10/0,693)$\cdot$0,024 $\approx$ 4,9$\times$10\^{}8 anos (cerca de 490 milhões de anos; com ln mais preciso, $\approx$ 4,8$\times$10\^{}8) & t/t$\frac12$ = 0,03463 $\Rightarrow$ t $\approx$ 4,88$\cdot$10\^{}8 anos & equivalente & ambos \\
\bottomrule\end{longtable}%
\begin{itemize}%
\item%
\textbf{Idade da amostra (relógio 208Pb/232Th))} Os dois usam o mesmo modelo: Pb/Th = e\^{}\{$\lambda$t\} - 1, logo t = t$\frac12$$\cdot$ln(1+R)/ln2. Pelos dados da prova, ln(1,024) = 7ln2 - 3ln5 = 7(0,693) - 3(1,609) = 0,024, que também é a aproximação ln(1+x) $\approx$ x. Assim, t/t$\frac12$ = 0,024/0,693 = 0,03463 e t = 0,03463 $\times$ 1,41$\times$10\^{}10 $\approx$ 4,883$\times$10\^{}8 anos, exatamente o valor do gabarito. A resolução chega ao mesmo número (4,883$\times$10\^{}8, conferido em Python) e o arredonda para 4,9$\times$10\^{}8 anos com 2 algarismos significativos. A diferença é só de arredondamento. Com ln(1,024) exato (0,02372), sai t $\approx$ 4,82$\times$10\^{}8 anos. A resolução cita esse valor como análise de sensibilidade, sem trocar a resposta principal. A razão t/t$\frac12$ $\approx$ 0,035 da resolução é apenas o arredondamento de 0,0346.%
\end{itemize}%
\paragraph{Pontos de dúvida}%
\begin{itemize}%
\item%
\textit{(solution)} A resposta principal foi arredondada para 4,9$\times$10\^{}8 anos (2 algarismos significativos), e o gabarito dá 4,88$\times$10\^{}8. O valor calculado é o mesmo (4,883$\times$10\^{}8). Um corretor rigoroso pode preferir 3 algarismos significativos, mas a diferença é só de arredondamento.%
\item%
\textit{(key)} O gabarito e a resolução tomam ln(1,024) = 0,024, que é o valor obtido com os ln 2 e ln 5 fornecidos (7ln2 - 3ln5). O valor exato é 0,02372, o que dá t $\approx$ 4,82$\times$10\^{}8 anos. Resposta na faixa de 4,8 a 4,9$\times$10\^{}8 anos deve ser aceita.%
\item%
\textit{(statement)} Os teores de ThO2 (9\%) e U3O8 (0,35\%) e a meia{-}vida do 238U não entram no cálculo, porque a razão Pb/Th dada já é molar. Esses dados funcionam como distratores. A resolução diz isso explicitamente e usa o U só numa checagem opcional (206Pb/238U $\approx$ 0,079).%
\end{itemize}%
\end{document}